\chapter{Introduzione}

\section*{Come usare questo manuale}

La documentazione di $\pi$Tantum \`e suddivisa in due parti, per
adattarsi a pubblici diversi:

\begin{itemize}
\item Il \textbf{Manuale generale} (Parte I, file
  \texttt{manual\_generale.pdf}) \`e pensato per chi usa
  $\pi$Tantum come strumento di lavoro: dirigenti scolastici,
  coordinatori d'orario, docenti che devono modificare un
  vincolo o capire perch\'e un orario \`e venuto in un certo
  modo. Tono discorsivo, casi d'uso concreti, niente jargon.
\item Le \textbf{Appendici tecniche} (Parte II, file
  \texttt{appendici\_tecniche.pdf}) sono per sviluppatori,
  amministratori IT e curiosi: architettura interna, modello
  dati, API REST, grammatica formale del DSL, dettaglio delle
  strategie di ottimizzazione e della diagnostica statistica.
\item Il \textbf{Manuale completo} (file
  \texttt{manual\_completo.pdf}) raccoglie entrambe le parti
  in un unico volume, utile per chi vuole stamparlo o averlo
  tutto a portata di mano.
\end{itemize}

Tutti e tre i PDF sono ricostruibili dai sorgenti
\texttt{docs/*.tex} con lo script \texttt{docs/build\_manual.sh}
(o \texttt{.bat} su Windows).

\section*{Cos'\`e $\pi$Tantum}

\emph{$\pi$Tantum} (alias \emph{Tempus Tantum}, codename interno:
\texttt{timetable}) \`e un sistema di generazione e gestione
dell'orario scolastico per un Liceo italiano. Il nome richiama il
celebre incipit dell'epistola I di Seneca a Lucilio:
``Omnia, Lucili, aliena sunt, tempus tantum nostrum est'' (Seneca,
\emph{Epistulae morales ad Lucilium}, I, 1) -- ``tutto ci appartiene
di altri, soltanto il tempo \`e nostro''. Il gioco grafico fra la
lettera greca $\pi$ e le due ``T'' di \emph{Tempus Tantum} riassume
l'idea: organizzare il tempo della scuola, dato che del tempo non
si dispone se non gestendolo.

\`E composto di tre layer:

\begin{enumerate}
  \item \textbf{Solver} (\texttt{experiments/}): pipeline CP-SAT
        (Google OR-Tools) con decomposizione spettrale per istanze
        grandi e cascata metaeuristica (LNS, SA, TS, ILS).
  \item \textbf{Web UI} (\texttt{webui/}): backend FastAPI su
        SQLite + frontend SvelteKit. \`E il punto di accesso quotidiano:
        gestisce CRUD, ottimizzazione, visualizzazione, drag-and-drop,
        assenze e supplenze.
  \item \textbf{Legacy} (\texttt{schedule/}): script monolitici da cui
        il solver \`e nato; mantenuti per riferimento storico.
\end{enumerate}

Il manuale copre tutti e tre i layer ma si concentra sulla web UI,
che ha assorbito quasi ogni funzionalit\`a operativa.

\section{Avvio rapido}

\paragraph{Windows.} Doppio click su \texttt{webui/start.bat}.
Apre due finestre cmd: backend \emph{uvicorn} su porta 8000, frontend
\emph{vite} su porta 5173.

\paragraph{Linux / macOS.}
\begin{lstlisting}[style=shell]
cd timetable/webui
./start.sh           # background; URLs sulla console
./stop.sh            # ferma backend + frontend
\end{lstlisting}

Nel browser, aprire \texttt{http://127.0.0.1:5173} e dalla
\textbf{Dashboard} importare un profilo g\`ia calcolato
(\texttt{small} per cominciare).

\section{Stato attuale del progetto}

Le funzionalit\`a principali, al momento di questa redazione:

\begin{itemize}
  \item 16 sezioni (tab) della UI, $\sim$118 endpoint REST.
  \item $\sim$30 tabelle SQLAlchemy con migrazioni idempotenti.
  \item 5 stati di vincoli (HARD/SOFT/PREFERITO/ENFORCED/ALLOWED).
  \item Vincoli logici DNF con DSL esteso (predicate atoms
        \texttt{aula:}, \texttt{materia:}, \texttt{classe:},
        \texttt{gruppo:}, \texttt{docente:}).
  \item Conflict detection automatico nel tab \emph{Vincoli}.
  \item Drag-and-drop con preview live (rosso/giallo/verde) e
        \emph{room-follows-lesson}.
  \item Gestione assenze settimanali con assegnazione supplenti
        drag-drop.
  \item Tab \emph{Monitor} con espansione lezione-per-lezione e
        slot picker matriciale.
\end{itemize}

% =========== Cap 2: Installazione ===========
